\documentclass{article}
\usepackage{tcee-repaso}
\tcetema[La OMC y el GATT]{3.B.33}{La OMC. Antecedentes y Organización actual. El GATT y los Acuerdos sobre el comercio de mercancías. Situación actual.}
\tcefecha{19/11/2026}
\begin{document}
\cabecerarepaso
\begin{columnas}

% ==========================================================================
\bloque[2,5]{Introducción}
\begin{itemize}
\item \textbf{Enganche}: «Hecha la ley, hecha la trampa».
  \begin{itemize}
  \item El GATT nació con la NMF y unas pocas excepciones, entre ellas la de seguridad (art.~XXI).
  \item 2025: EE.~UU. pone aranceles a casi todos e invoca la seguridad nacional; comercio NMF 83\,\% (2022) $\to$ 72\,\% (feb. 2026).
  \end{itemize}
\item \textbf{Relevancia}: 166~miembros, $\approx$98\,\% del comercio; aranceles industriales del 40\,\% a $<$4\,\%.
  \begin{itemize}
  \item España: GATT 1963 \vertema{4.A.27}; negocia a través de la UE \vertema{3.B.41}; aceitunas negras (DS577).
  \end{itemize}
\item \textbf{Contextualización}: \concepto{OMC}: Marrakech, en vigor 1-1-1995, Ginebra.
  \begin{itemize}
  \item Bretton Woods: FMI \vertema{3.B.22}, BIRF \vertema{3.B.30} y la OIC fallida; quedó el \concepto{GATT} (1947), provisional 47~años.
  \item La OMC añade institución y \emph{todo único}, consenso negativo en las diferencias y más ámbito (servicios, ADPIC \vertema{3.B.34}).
  \end{itemize}
\item \textbf{Problemática}: ¿por qué un acuerdo si el libre comercio ya conviene? ¿Cómo decide? ¿Qué reglas y excepciones? ¿Por qué está atascada?
\item \textbf{Estructura}: \hyperref[rep:I]{I} institución; \hyperref[rep:II]{II} reglas; \hyperref[rep:III]{III} crisis.
\end{itemize}
\esquemacante{%
\apartadocante{Introducción}{2,5}
\apartadocante{I. La OMC: antecedentes y organización}{8,5}
\subcante{1. Antecedentes de la OMC}{4}
\subcante{2. Organización actual}{4,5}
\apartadocante{II. El GATT y los acuerdos sobre mercancías}{10}
\subcante{1. Acuerdos multilaterales}{8}
\subcante{2. Acuerdos plurilaterales}{2}
\apartadocante{III. Negociaciones y situación actual}{8}
\subcante{1. Evolución de las negociaciones}{2,5}
\subcante{2. Situación actual}{5,5}
\apartadocante{Conclusión}{1}}

% ==========================================================================
\bloque[8,5]{I. La OMC: antecedentes y organización\label{rep:I}}
\begin{ideaclave}
La OMC resuelve un dilema del prisionero: cada país grande protege a costa de los demás y todos pierden. El GATT lo resolvió con reciprocidad y no discriminación; la OMC añadió un juez.
\end{ideaclave}

\sub{I.1. Antecedentes de la OMC\label{rep:I.1}}
\subsub{a) Del bilateralismo a la guerra arancelaria}
\begin{itemize}
\item \emph{Siglo XIX}: tratados bilaterales con \concepto{NMF}.
  \begin{itemize}
  \item \autor{Cobden}-\autor{Chevalier} (1860): NMF incondicional que generaliza rebajas. Proteccionismo desde 1880.
  \end{itemize}
\item \emph{Años 30}: \concepto{Ley \autor{Smoot}-\autor{Hawley}} (1930), $\approx$20.000~productos, represalias.
  \begin{itemize}
  \item Comercio $-$2/3 en valor (1929-1933).
  \end{itemize}
\end{itemize}
\subsub{b) ¿Por qué un acuerdo? \autor{Bagwell} y \autor{Staiger} (1999)}
\begin{itemize}
\item \emph{Idea}: un país grande gana con un arancel mejorando su RRI \vertema{3.B.7}.
  \begin{itemize}
  \item Si todos lo hacen, pierden (\autor{Johnson}, 1953) \vertema{3.B.8}.
  \item El acuerdo solo tiene que corregir esa \emph{externalidad de la RRI}.
  \end{itemize}
\item \emph{Resultado}: el arancel de Nash supera al \concepto{políticamente óptimo} (el que se elegiría ignorando la RRI).
  \begin{itemize}
  \item Los 2~gobiernos ganan con el óptimo político.
  \item Dilema del prisionero: (5, 5), (0, 6), (6, 0), \colorbox{tceesalvia}{(1, 1)}.
  \end{itemize}
\item \emph{Implicaciones}:
  \begin{itemize}
  \item \emph{reciprocidad}: rebajas que dejan la RRI igual;
  \item \emph{NMF}: precio mundial único;
  \item autoejecutable: represalia creíble y acotada $\Rightarrow$ importa el juez.
  \end{itemize}
\item \emph{Extensiones y evidencia}: compromiso (\autor{Maggi}-\autor{Rodríguez-Clare}, 1998); \autor{Broda}-\autor{Limão}-\autor{Weinstein} (2008).
\item \emph{Valoración}: explica la arquitectura del GATT, no los temas nuevos.
\end{itemize}
\subsub{c) Orígenes del GATT: de La Habana a Ginebra}
\begin{itemize}
\item Carta de La Habana (1948): comercio + empleo; EE.~UU. no la ratifica (1950).
\item \concepto{GATT}: 30-10-1947, 23~países, 45.000~concesiones.
  \begin{itemize}
  \item Aplicado de forma provisional desde el 1-1-1948; nunca entró en vigor como tratado.
  \end{itemize}
\end{itemize}
\subsub{d) Las rondas de negociación}
\begin{itemize}
\item \emph{1.ª etapa} (1947-1961): 5~rondas \emph{producto a producto}, proveedor principal.
  \begin{itemize}
  \item Ginebra $-$32\,\%; después, 2-3,5\,\%; crestas. España, 1963.
  \end{itemize}
\item \emph{2.ª etapa}: fórmulas y barreras no arancelarias.
  \begin{itemize}
  \item \emph{Kennedy} (1964-1967, 62~países): lineal, $-$36\,\%; Código Antidumping; Parte IV.
  \item \emph{Tokio} (1973-1979, 102): fórmula suiza, $-$1/3 hasta 4,7\,\%; códigos.
  \item Neoproteccionismo: Multifibras, RVE.
  \end{itemize}
\end{itemize}
\graficorepaso[0.75\linewidth]{formula-suiza}{Fórmula suiza $t_1=at_0/(a+t_0)$: $a$ es el techo; recorta más los altos}
\begin{itemize}
\item \emph{3.ª etapa}: \emph{Uruguay} (1986-1994, 123~países).
  \begin{itemize}
  \item $-$40\,\% (6,3 $\to$ 3,8\,\%); consolidación 78 $\to$ 99\,\% (PD) y 21 $\to$ 73\,\% (PED).
  \item Agricultura y textiles; AGCS, ADPIC, MIC; Marrakech (15-4-1994).
  \end{itemize}
\end{itemize}
\graficorepaso[0.92\linewidth]{rondas-gatt}{Recorte arancelario medio por ronda y participantes: de 23~a 123~países}

\sub{I.2. Organización actual\label{rep:I.2}}
\subsub{a) Marco jurídico}
\begin{itemize}
\item Acuerdo de Marrakech (16~arts.) + anexos:
  \begin{itemize}
  \item 1A mercancías, 1B AGCS, 1C ADPIC; 2~ESD; 3~MEPC: multilaterales (\emph{todo único});
  \item 4~plurilaterales.
  \end{itemize}
\item GATT 1994 $\neq$ GATT 1947 (art.~II.4).
\end{itemize}
\graficorepaso[0.92\linewidth]{acuerdos-omc}{El anexo 1A (mercancías) frente a servicios, ADPIC, ESD, MEPC y plurilaterales}
\subsub{b) Marco institucional}
\subsubsub{Características}
\begin{itemize}
\item Funciones (art.~III); 166~miembros (Comoras y Timor Oriental, 2024; también territorios aduaneros).
\item Un voto, \concepto{consenso}; cuotas según comercio.
\item Presupuesto $\approx$205~M CHF, $\approx$620~personas; EE.~UU. con atrasos.
\item DG \autor{Ngozi Okonjo-Iweala}, 2.º mandato desde 1-9-2025.
\end{itemize}
\subsubsub{Estructura}
\begin{itemize}
\item Conferencia Ministerial (cada 2~años) $\to$ Consejo General (también OSD y OEPC) $\to$ Consejos de Mercancías, Servicios y ADPIC $\to$ comités.
\item Secretaría sin poder de decisión.
\end{itemize}
\graficorepaso[0.92\linewidth]{estructura-omc}{3 niveles: Conferencia, Consejo General (OSD, OEPC), consejos y comités}
\subsubsub{La solución de diferencias}
\begin{itemize}
\item Consultas (60~días) $\to$ grupo especial (3-5 miembros, 6-9 meses) $\to$ apelación (7~miembros, 60-90 días) $\to$ plazo prudencial $\to$ compensación o represalias.
\item \concepto{Consenso negativo}: casi automático. 646~casos (1995-2026).
\end{itemize}
\graficorepaso[0.92\linewidth]{procedimiento-diferencias}{Consultas, grupo especial, apelación (bloqueada: MPIA), cumplimiento y represalias}
\subsubsub{Examen de políticas y transparencia}
\begin{itemize}
\item Notificaciones; exámenes cada 3~años (4~primeros), 5 (16~siguientes), 7 (resto).
\item Vigilancia del G-20 (nov. 2025): nuevas restricciones, 14,3\,\% de sus importaciones.
\end{itemize}

% ==========================================================================
\bloque[10]{II. El GATT y los acuerdos sobre mercancías\label{rep:II}}
\begin{ideaclave}
Principios con válvulas (integración, desarrollo, salvaguardias, defensa comercial, seguridad): el acuerdo aguanta mientras las válvulas se usen con mesura.
\end{ideaclave}

\sub{II.1. Acuerdos multilaterales\label{rep:II.1}}
\subsub{a) El GATT: estructura y principios}
\begin{itemize}
\item \emph{Estructura}: 38~arts. en 4~partes (NMF y listas; reglas; integración y adhesión; desarrollo, 1965).
\item \emph{Principios}:
  \begin{itemize}
  \item \concepto{no discriminación}: NMF (art.~I) y trato nacional (art.~III);
  \item \concepto{consolidación} (art.~II); \concepto{transparencia} (art.~X);
  \item \concepto{prohibición de cuotas} (art.~XI; los contingentes arancelarios sí);
  \item \concepto{liberalización progresiva}.
  \end{itemize}
\item \emph{Excepciones}:
  \begin{itemize}
  \item art.~XXIV: integración ($\approx$90\,\% del comercio; no subir la protección; compensar: ajos de Croacia); \autor{Kemp}-\autor{Wan} (1976) \vertema{3.B.10};
  \item \concepto{cláusula de habilitación} (1979; SPG \vertema{3.B.41}); \emph{waivers};
  \item balanza de pagos (arts.~XII y XVIII.B: EE.~UU., sección~122, 2026);
  \item salvaguardias (XIX); generales (XX: MAFC \vertema{3.B.31});
  \item \concepto{seguridad} (XXI): \emph{Rusia — Tránsito} (2019), no del todo autojuzgable.
  \end{itemize}
\end{itemize}
\subsub{b) Acuerdos sobre cuestiones y sectores específicos}
\graficorepaso[0.95\linewidth]{acuerdos-especificos}{Los 13~acuerdos específicos del anexo 1A en 5~grupos, con su año y su idea clave}
\subsubsub{Procedimientos aduaneros}
\begin{itemize}
\item Valoración (valor de transacción); origen (última transformación sustancial; cuenco de espaguetis); licencias; inspección previa.
\item \concepto{AFC}: Bali 2013, en vigor 22-2-2017; 89,6\,\% aplicado; costes $-$13-15\,\%.
\end{itemize}
\subsubsub{Reglamentos técnicos y sanitarios}
\begin{itemize}
\item MSF: ciencia, \emph{Codex}; hormonas (DS26). OTC: normas internacionales.
\item CM14: trato especial MSF/OTC.
\end{itemize}
\subsubsub{Defensa comercial}
\begin{itemize}
\item \concepto{Antidumping}: margen $\geq$2\,\% del precio de exportación, daño, causalidad.
  \begin{itemize}
  \item Apoyo $\geq$25\,\% de la producción y más que los opuestos; importaciones \emph{de minimis} 3\,\%; 5~años.
  \end{itemize}
\item \concepto{SMC}: solo las subvenciones específicas.
  \begin{itemize}
  \item \emph{Prohibidas}: a la exportación y de contenido local.
  \item \emph{Recurribles}: impugnables si causan efectos desfavorables (daño, anulación o menoscabo, perjuicio grave).
  \item \emph{No recurribles}: hasta el 31-12-1999. Vías: OSD o derechos compensatorios.
  \end{itemize}
\item \concepto{Salvaguardias}: daño grave, NMF, 4 + 4~años, compensación; sin RVE.
\end{itemize}
\graficorepaso[0.92\linewidth]{defensa-comercial}{Investigaciones: 7.471~antidumping, 816~compensatorias, 455~salvaguardias (1995-2025)}
\subsubsub{Inversión y sectores}
\begin{itemize}
\item \emph{MIC}: prohíbe contenido local y equilibrio de divisas.
\item \emph{Textiles}: Multifibras $\to$ ATV; cuotas fuera el 1-1-2005.
\item \emph{Pesca} (CM12, 2022): en vigor 15-9-2025 ($>$120~aceptaciones).
  \begin{itemize}
  \item Prohíbe ayudas a la pesca INDNR, a poblaciones sobreexplotadas y en alta mar.
  \item «Segunda ola» pendiente; caduca en 2029 sin ella. España: mayor flota de la UE.
  \end{itemize}
\end{itemize}
\subsubsub{El Acuerdo sobre la Agricultura}
\begin{itemize}
\item Arancelización y $-$36\,\% en 6~años (PD), $-$24\,\% en 10 (PED).
\item \emph{Cajas}:
  \begin{itemize}
  \item \textcolor{tceedorado}{ámbar}: MGA $-$20\,\%, \emph{de minimis} 5/10\,\%;
  \item \textcolor{tceeazul}{azul}, \textcolor{tceeverde}{verde} y de desarrollo.
  \end{itemize}
\item Nairobi: fin de las subvenciones a la exportación (PD ya; PED 2018/2022).
\item Existencias públicas: cláusula de paz (Bali), sin solución en la CM14.
\end{itemize}
\subsub{c) La economía de las salvaguardias (\autor{Bagwell}-\autor{Staiger}, 1990)}
\begin{itemize}
\item \emph{Idea}: el acuerdo debe ser autoejecutable.
  \begin{itemize}
  \item En las oleadas de importaciones la tentación de desviarse crece y puede superar el coste de la represalia.
  \end{itemize}
\item \emph{Resultado}: permitir entonces una protección limitada y temporal, por debajo del arancel de Nash.
  \begin{itemize}
  \item Es la \concepto{cláusula de escape} (art.~XIX): sin válvula, el acuerdo salta.
  \end{itemize}
\item \emph{Extensiones}: usarla debe tener precio (\autor{Rosendorff}-\autor{Milner}, 2001); antidumping como salvaguardia encubierta (\autor{Bown}-\autor{Crowley}, 2013).
\item \emph{Valoración}: la excepción de seguridad no tiene precio ni plazo.
\end{itemize}
\subsub{d) Listas de compromisos}
\begin{itemize}
\item $\approx$22.500~páginas (Ronda Uruguay).
\item \concepto{Excedente de consolidación} (consolidado/aplicado):
  \begin{itemize}
  \item UE 5/5\,\%; EE.~UU. 3,4/3,3;
  \item India 50,8/16,2 (74\,\% consolidado); Brasil 31,4/12.
  \end{itemize}
\item Mundo: NMF simple 13,1\,\% (1996) $\to$ 8,9\,\% (2024).
\end{itemize}
\graficorepaso[0.92\linewidth]{aranceles-consolidados}{Consolidado frente a aplicado (2024): margen grande en los emergentes}

\sub{II.2. Acuerdos plurilaterales\label{rep:II.2}}
\begin{itemize}
\item \emph{Qué son}: solo obligan a quien los firma; la vía cuando el consenso es imposible.
\item \emph{Anexo 4} (reglas propias, solo para sus partes):
  \begin{itemize}
  \item Aeronaves Civiles (1979): sin aranceles en aeronaves y piezas; 33~signatarios;
  \item Contratación Pública (ACP): compras de las administraciones; 22~partes, 49~miembros \vertema{3.B.34};
  \item bovino y lácteos: fin en 1997.
  \end{itemize}
\item \emph{Cero por cero} (extendidos por NMF):
  \begin{itemize}
  \item \concepto{ATI} (1996; 84~participantes, 97\,\%; ampliado en 2015 a 201~productos);
  \item farmacéuticos, químicos; Bienes Ambientales (suspendido);
  \item problema: el \emph{polizón}: quien no participa se beneficia sin abrir;
  \item solo compensan con \emph{masa crítica}: que los participantes hagan casi todo el comercio del sector.
  \end{itemize}
\item \emph{Nuevos}:
  \begin{itemize}
  \item Facilitación de Inversiones (131): la apoyan 165~de 166, pero sin consenso para el anexo 4 (la India, según la prensa);
  \item Comercio Electrónico (66~en Yaundé, 68~hoy; aplicación provisional con 45~aceptaciones).
  \end{itemize}
\end{itemize}

% ==========================================================================
\bloque[8]{III. Negociaciones y situación actual\label{rep:III}}
\begin{ideaclave}
3~acuerdos multilaterales nuevos en 30~años, sin juez de apelación desde 2019 y con su impulsor fuera; aun así, casi 3/4 del comercio sigue sus reglas.
\end{ideaclave}

\sub{III.1. Evolución de las negociaciones\label{rep:III.1}}
\begin{itemize}
\item \emph{Ronda de Doha} (2001), «del Desarrollo»: todo único, a cerrar en 2005.
  \begin{itemize}
  \item Cancún (2003): fracasa por la agricultura y los temas de Singapur (abandonados en 2004).
  \item Julio de 2008: casi se cierra; se rompe por la salvaguardia agrícola especial (India y China frente a EE.~UU.).
  \item Después, solo acuerdos parciales; en Nairobi ya no se reafirma el mandato: para muchos, muerta.
  \item Fondo: rebajar aranceles ya bajos rinde poco.
  \end{itemize}
\item \emph{Conferencias Ministeriales} (cada 2~años):
  \begin{itemize}
  \item Singapur 1996 (temas de Singapur; ATI); Ginebra 1998 (moratoria digital); Seattle 1999 (fracaso);
  \item \concepto{Doha} 2001: ronda del Desarrollo, todo único; entra China;
  \item Cancún 2003 (fracaso); Hong Kong 2005; julio de 2008 (salvaguardia agrícola especial);
  \item \textbf{Bali 2013} (AFC); \textbf{Nairobi 2015} (subvenciones agrícolas a la exportación; ATI ampliado);
  \item Buenos Aires 2017; \textbf{Ginebra 2022} (pesca; vacunas); Abu Dabi 2024.
  \end{itemize}
\end{itemize}
\subsub{La CM14 (Yaundé, 2026) y lo que vino después}
\begin{itemize}
\item 26-30 de marzo: 3~decisiones (pequeñas economías; MSF/OTC; seguir con la pesca).
\item Paquete de Yaundé (reforma, comercio electrónico, PMA), al Consejo General.
\item Expiran las moratorias digital (desde 1998; Brasil y Turquía se oponen, según EE.~UU.) y de no infracción de los ADPIC.
  \begin{itemize}
  \item En contra (India, Indonesia, Sudáfrica): recaudación y espacio para la política digital.
  \item A favor: poca recaudación; encarece insumos digitales.
  \item 19~miembros sin aranceles digitales entre sí (mayo).
  \end{itemize}
\item Reforma: 4~líneas (principios, decisiones, desarrollo, igualdad de condiciones); CM15 sin sede (Riad, 2028).
\end{itemize}

\sub{III.2. Situación actual\label{rep:III.2}}
\subsub{a) La parálisis negociadora}
\begin{itemize}
\item Consenso (166~vetos) y todo único: cada tema, atado a los demás.
\item Asimetría: en el GATT los PD negociaban y extendían las ventajas a los PED sin pedir nada; hoy China e India son potencias.
  \begin{itemize}
  \item Trato especial autodeclarado: China pide lo mismo que Malí; renuncian Singapur, Corea, Brasil y China (sep. 2025).
  \end{itemize}
\item Intereses opuestos: los PED quieren la agricultura de los ricos; los ricos, servicios, propiedad intelectual y compras públicas.
\item Reglas viejas: nada eficaz sobre subvenciones industriales, empresas públicas y exceso de capacidad, datos, clima.
\end{itemize}
\subsub{b) La crisis de la solución de diferencias}
\begin{itemize}
\item EE.~UU. bloquea el OA desde 2017: sin quórum el 11-12-2019; vacío desde 30-11-2020.
  \begin{itemize}
  \item Quejas: se extralimita, incumple plazos, crea precedentes.
  \end{itemize}
  \begin{itemize}
  \item Apelación «al vacío»: $\approx$30~pendientes (DS623, 2026).
  \item Demandas: 20-40/año $\to$ 5-13.
  \end{itemize}
\item \concepto{MPIA} (2020, art.~25~ESD): 62~miembros, 60\,\% del comercio (Camboya, 2026).
  \begin{itemize}
  \item Sin EE.~UU. ni India; 1.\textsuperscript{er} laudo DS611 (2025).
  \end{itemize}
\end{itemize}
\graficorepaso[0.92\linewidth]{diferencias-anio}{Solicitudes de consultas por año: máximo en 1997 (50); desde el bloqueo, 5-13}
\subsub{c) El desafío unilateral: aranceles de EE.~UU. (2025-2026)}
\begin{itemize}
\item IEEPA 2025: recíprocos (techo UE 15\,\%); el Supremo los anula (20-2-2026).
\item Sustitutos: sección~122 (10\,\%, feb.-jul.); sección~301 (10/12,5\,\%, desde 24-7-2026); 232 (acero, aluminio, cobre, autos).
\item \emph{Encaje}: violan los arts.~I-II; no son art.~XXIV; EE.~UU. invoca el XXI (DS633, DS634\ldots); sin juez.
\item \emph{Respuestas}: China, represalias; UE, Japón, RU, techos. Comercio +4,2\,\% (2025), pero desviado.
\item \emph{Lectura del modelo}: vuelve el país grande que usa su poder de mercado; NMF 83 $\to$ 72\,\% \vertema{3.B.21}.
\end{itemize}
\subsub{d) Multilateralismo frente a regionalismo}
\begin{itemize}
\item $\approx$390~ACR notificados; todos los miembros tienen alguno \vertema{3.B.35}.
\item ¿\emph{Building} o \emph{stumbling blocks}? (\autor{Bhagwati}):
  \begin{itemize}
  \item estorban: desvían comercio \vertema{3.B.10}, crean intereses en las preferencias, normas de origen;
  \item construyen: avanzan entre pocos en lo que la OMC no puede; \emph{regionalismo abierto}.
  \end{itemize}
\item Solo el 17\,\% del comercio usa preferencias: la NMF sobrevivió al regionalismo; la erosiona el unilateralismo.
\item Las CGV piden reglas profundas \vertema{3.B.18}.
\item \emph{UE y España}: DS577 aceitunas (13,64~M\,\$/año, 2025); DS630 vehículos eléctricos; DS593 palma; DS639 MAFC.
\end{itemize}

% ==========================================================================
\bloque[1]{Conclusión}
\begin{itemize}
\item \textbf{Recapitulación}:
  \begin{itemize}
  \item corrige un fallo de coordinación;
  \item reciprocidad, NMF y consolidación: 8~rondas llevan los aranceles industriales casi a cero;
  \item la OMC le da organización, más ámbito y un juez;
  \item principios con válvulas.
  \end{itemize}
\item \textbf{Relevancia}: para una economía mediana como España, en la UE, las reglas valen más que para una potencia: sin ellas, solo cuenta el tamaño.
\item \textbf{Extensiones}: \vertema{3.B.34} \vertema{3.B.7} \vertema{3.B.10} \vertema{3.B.35} \vertema{3.B.41} \vertema{4.A.24}.
\item \textbf{Opinión}: la ventaja comparativa, «verdadera y no trivial» (\autor{Samuelson}, 1969); la OMC no se sustituirá, pero no volverá a ser la de 1995.
  \begin{itemize}
  \item OMC de geometría variable: plurilaterales abiertos, MPIA y límites a la excepción de seguridad.
  \end{itemize}
\item \textbf{Idea final}: el GATT no eliminó las trampas: les puso precio y plazo.
  \begin{itemize}
  \item Mientras la mayoría prefiera eso a la ley del más fuerte, las reglas seguirán sirviendo.
  \end{itemize}
\end{itemize}

% ==========================================================================
\anexorepaso
\bloque{Datos clave}
\begin{itemize}
\item \dato{166 miembros}{OMC, 2026}; 22~adhesiones en curso.
\item \dato{646 diferencias}{1995-2026}; 2025: 13. MPIA: \dato{62 miembros, 60\,\% del comercio}{2026}.
\item NMF aplicado: \dato{UE 4,9\,\%; EE.~UU. 3,3; China 7,5; India 16,0}{OMC, 2025}.
\item Comercio en NMF: \dato{83\,\% (2022) $\to$ 72\,\% (2026)}{OMC}.
\item G-20: nuevas restricciones \dato{14,3\,\% de sus importaciones}{nov. 2025}.
\end{itemize}
\bloque{Autores y fechas}
\hito{1860}{\autor{Cobden}-\autor{Chevalier}: NMF incondicional}
\hito{1930}{\autor{Smoot}-\autor{Hawley}}
\hito{1947}{GATT (23~países)}
\hito{1976}{\autor{Kemp}-\autor{Wan}: uniones aduaneras}
\hito{1995}{OMC (Marrakech, 1994)}
\hito{1999}{\autor{Bagwell}-\autor{Staiger}: externalidad de la RRI}
\hito{2019}{Bloqueo del Órgano de Apelación}
\hito{2026}{CM14 (Yaundé); fin de la moratoria digital}
\bloque{Preguntas del tribunal}
\begin{enumerate}
\item \textbf{¿Por qué un acuerdo si el libre comercio ya conviene?} Externalidad de la RRI (\autor{Bagwell}-\autor{Staiger}) o compromiso.
\item \textbf{¿Consolidado o aplicado?} Techo frente a lo que se cobra: excedente de consolidación.
\item \textbf{¿Es legal un arancel por seguridad nacional?} Art.~XXI; no del todo autojuzgable; sin OA, no hay sentencia firme.
\item \textbf{¿Antidumping o salvaguardia?} Desleal, daño importante, selectivo; frente a aumento de importaciones, daño grave, NMF.
\item \textbf{¿Por qué fracasó Doha?} Todo único, consenso, asimetría, agricultura.
\item \textbf{¿Qué es un plurilateral «cero por cero»?} Rebaja sectorial entre unos pocos extendida por NMF (ATI): exige masa crítica.
\end{enumerate}
\begin{trampa}
\begin{itemize}
\item 23~fundadores del GATT; España \emph{no}: 1963.
\item Anexo 4~hoy: \emph{2} plurilaterales; el ATI \emph{no} es del anexo 4.
\item GATT aplicado de forma provisional desde el 1-1-1948; OMC, 1-1-1995.
\item Fórmula suiza: $a$ es el arancel \emph{máximo} final.
\item Salvaguardias: máximo \emph{8} años; prohibidas las RVE.
\item OA sin quórum el \emph{11-12-2019}; MPIA: \emph{art.~25} del ESD.
\end{itemize}
\end{trampa}

\end{columnas}
\end{document}
